% ─────────────────────────────────────────────────────────────────────────────
\section{Experimental Setup}
\label{sec:setup}

\paragraph{Datasets and roles.}
\emph{UAV-VisLoc}~\cite{xu2024uavvisloc} (11 flights over China, 6{,}742 nadir
images, 11 geo-referenced satellite maps of about 0.3\,m/px, altitude and
heading metadata) is the main benchmark for continuous-map localization and all
mechanistic analyses. \emph{World-UAV}~\cite{wu2025uavgeoloc}, Rot split
(one region, 3 heights $\times$ 72 headings; satellite tiles mosaicked into a
continuous map), stresses rotation and scale with an \emph{unknown} heading.
\emph{AnyVisLoc}~\cite{ye2026anyvisloc} (low-altitude, multi-view, aerial and
satellite references) probes the viewpoint validity envelope under its
known-attitude setting. Parameters are selected on UAV-VisLoc only and frozen
for the other two datasets.

\paragraph{Protocol (UAV-VisLoc).}
For each site, 100 queries are sampled uniformly along the flight. The
altitude-to-GSD factor $k$ (metadata lack intrinsics) is estimated once on a
calibration site (site 05) that is excluded from the reported results. The
main setting simulates an INS/VO prior: the prior centre is the true position
plus an offset drawn uniformly in a disk of radius $R/2$, and the search window
has radius $R=1$\,km. Priors are drawn with a fixed per-query seed, so every
method sees identical priors (paired comparison). Queries are stratified by
the number of buildings \emph{observed in the query} (no ground truth is used
for selection).

\paragraph{Reference points and baselines.}
(i) \emph{Random}: the prior centre, and a uniform random point in the window,
on the same queries. (ii) \emph{M0}: the conference descriptor-retrieval
pipeline (learned segmentation, CDT, MFCA, $k$-NN voting), restricted to the
same prior window. (iii) \emph{Oracle structure}: the query line structure
replaced by the map's own line structure under the true footprint, isolating
image$\rightarrow$structure errors (UAV vs.\ satellite line extraction) from
registration errors. (v) \emph{Ablations}: isotropic chamfer (no orientation
channels), no graph verification, and graph verification without the Ekeland
term. (iv) \emph{Local
sanity}: search restricted to 50\,m around the ground truth, testing whether
the objective peaks at the correct pose at all.

\paragraph{Metrics.}
Median, mean and 95th-percentile horizontal error; Success@25/50/100\,m;
candidate recall (a top-5 peak within 25\,m of the truth); integrity AUROC,
ECE and AURC. Confidence intervals are 95\% cluster bootstrap over sites where
applicable; paired differences use the same queries.

\paragraph{Hardware.}
One NVIDIA A100 40\,GB (shared; about 20\,GB available), 16 CPU cores.

% ─────────────────────────────────────────────────────────────────────────────
\section{Results}
\label{sec:exp_main}
\TBD{G0 gate first (UAV-VisLoc sites 01 and 11). Filled only from results/reg/*.csv.}

\subsection{Why per-element matching fails (M0)}
\label{sec:why_m0}
The conference pipeline (M0) retrieves individual CDT triangles (5-D
angle/Ekeland descriptors, $\ell_1$ $k$-NN) and votes for a satellite patch.
On five terrain-diverse UAV-VisLoc sites (100 queries each) it fails
everywhere: the fraction of queries within 500\,m never exceeds $3.1\%$ and the
median error is $3.3$--$4.4$\,km, although the true location is covered by
mapped buildings in $93$--$99\%$ of frames (full tables in the supplement).
A truth-anchored diagnosis on the Changjiang-20 map
($\approx4.8\times10^5$ building triangles) shows why: the geometrically
correct satellite element is in the top-20 for $0\%$ of query triangles, and
enriching the descriptor with size, a log-polar building-constellation
context, building-level matching or per-dimension standardization shrinks the
true-match distance by up to $16\times$ yet leaves retrieval at $0\%$. The
descriptor is too low-entropy to be \emph{unique} among $10^5$--$10^6$
candidates: impostors approach the query as fast as the true match does.
Restricting M0 to a small prior window ($\approx0.2$\,km$^2$) bounds the error
mechanically (146\,m) but provides almost no identification (2.1\% within
30\,m vs.\ 1.4\% for a uniform guess), and integrating the per-frame votes over a
trajectory with a particle filter does not converge, because the true position
receives votes in only $12\%$ of frames. These results define the requirement
that motivates StructReg: the observation must use the \emph{joint}
configuration of the observed buildings.

\subsection{World-UAV: rotation and scale}
\label{sec:exp_worlduav}
\TBD{}

\subsection{AnyVisLoc: viewpoint validity envelope}
\label{sec:exp_anyvisloc}
\TBD{}

\section{Integrity and Efficiency}
\label{sec:efficiency}
\TBD{}
